\begin{frame}{Q-함수를 신경망으로 근사: 손실함수}
  \begin{itemize}
    \item DQN의 아이디어: 테이블 대신 \textbf{신경망}으로 $Q(s,a;\theta)$ 근사.
    \item 손실함수 = Week 6의 \textbf{TD 오차를 제곱한 것}:
  \end{itemize}
  \[
  J(\theta) = \mathbb{E}\left[\left(r + \gamma \max_{a'} Q(s',a';\theta)
  - Q(s,a;\theta)\right)^2\right]
  \]
  \begin{itemize}
    \item 경사하강법(역전파)으로 $\theta$ 갱신 --- \textbf{형태는}
      지도학습 회귀와 같음.
    \item 그러나 ``정답''인 목표값 $r + \gamma \max_{a'} Q(s',a';\theta)$ 자체가
      지금 학습 중인 $\theta$로 계산된다.
  \end{itemize}
\end{frame}
